\clearpage
\fancyhead[L]{\small BODE. ASTRONOMISCHES JAHRBUCH.}
\SourceLeaf{509}
\begin{center}\small Göttingische gelehrte Anzeigen. Stück 155. S. 1545..1552. 26. September 1808.\\[3pt]\rule{.72\linewidth}{.3pt}\end{center}

Beim Verfasser und in Kommission bei \textsc{Fr. Braunes}. \emph{Astronomisches Jahrbuch für das Jahr} 1810, \emph{nebst einer Sammlung der neuesten in die astronomischen Wissenschaften einschlagenden Abhandlungen und Nachrichten. Von} \textsc{J. E. Bode}. 1807. 268 Seiten im Oktavformat und eine Kupfertafel. Berlin.

Der erste Teil dieses geschätzten und für jeden Astronomen unentbehrlichen Jahrbuchs, der den astronomischen Kalender enthält, ist ebenso eingerichtet wie im vorangegangenen Jahrgang. Der verehrte Herausgeber hat sich jedoch seit jeher unermüdlich darum bemüht, es brauchbarer zu machen, und Vorschläge und Wünsche in dieser Richtung berücksichtigt. Deshalb erlauben wir uns, hier einige solche Wünsche vorzutragen, die, wie wir sicher wissen, auch mehrere andere Astronomen seit Langem hegen. Die ausführliche Erläuterung des Kalenders und seines Gebrauchs wurde seit dem Jahrgang 1800 nicht wieder abgedruckt. Wir finden das sehr lobenswert, denn dieser ziemlich umfangreiche Aufsatz würde Platz einnehmen, den man besser nutzen könnte. Außerdem darf man \emph{das Meiste}, was er enthält, ohnehin bei jedem als bekannt voraussetzen, der den Himmel beobachten will. Einiges aus dieser Erläuterung betrifft jedoch die willkürlichen und nicht selbstverständlichen Festlegungen des Kalenders. Das sollte man wirklich zu einer ständigen Rubrik machen oder auf andere Weise in den Kalender aufnehmen. Dass zum Beispiel alle Planetenpositionen für Mitternacht wahrer Berliner Zeit angesetzt sind, mit Ausnahme von Venus und Merkur, deren Positionen für Mittag gelten, errät niemand von selbst. Wer den Jahrgang 1800 nicht nachschlagen kann, bleibt darüber im Ungewissen. Ebenso steht es mit der Angabe, dass die Bedeckungen der Fixsterne durch den Mond für die wahre Zeit gelten und nicht für die mittlere, die heute bei fast allen Astronomen üblich ist. Bei der scheinbaren Schiefe der Ekliptik wird die Nutation stets mit umgekehrtem Vorzeichen angegeben. Ohne Erklärung ist das nicht selbstverständlich. Da die scheinbare Schiefe die Nutation enthält, würde man eher annehmen, dass man diese davon abziehen muss, um die mittlere Schiefe zu erhalten. Bei der Angabe im Jahrbuch verhält es sich aber umgekehrt. Auch können wir es nicht gutheißen, dass Herr \textsc{Bode} die Schiefe der Ekliptik noch immer nach den älteren Sonnentafeln von Herrn \textsc{von Zach} angibt. Nach den neueren Beobachtungen, die englische, französische und italienische Astronomen schon seit vielen Jahren anstellen, ist sie um $8''$ zu groß. Daher sind auch alle Deklinationen der Sonne im Jahrbuch zu groß und ergeben, wenn man sich auf sie verlässt, je nach Jahreszeit unrichtige Polhöhen. Es wäre auch zu wünschen, dass Herr \textsc{Bode} die unbequeme Zeitangabe nach Vormittags- und Nachmittagsstunden durch die unter Astronomen übliche ersetzt. Wir kennen mehrere Fälle, in denen Sternbedeckungen und Verfinsterungen der Jupitermonde allein wegen der Angabe in bürgerlicher Zeit verpasst wurden. Auf den Wunsch eines Rezensenten im 41. Stück der Literaturzeitung von 1797 hin hat Herr \textsc{Bode} seit 1800 bei den Planetenpositionen auch die heliozentrischen Längen und Breiten aufgenommen, dafür aber die Rektaszensionen weggelassen. Damit waren wir nie zufrieden. Wir glauben, dass die geozentrische Rektaszension und Deklination das Allerwesentlichste in einer astronomischen Ephemeride sind. Lieber würden wir alle anderen Spalten zu den Planeten entbehren als diese beiden notwendigen Größen. Gewiss würde jeder Astronom aus mancherlei Gründen begrüßen, wenn auch die Abstände von der Erde angegeben würden, die zur Bestimmung der Aberration unentbehrlich sind. Die Gründe, aus denen jener Rezensent die heliozentrischen Positionen wünschte, scheinen uns nicht so erheb\SourceLeaf{510}lich. Wir glauben, dass sich von diesen beiden Spalten wenig Gebrauch machen lässt, zumal Herr \textsc{Bode} unter den monatlichen Erscheinungen die Oppositionen und Konjunktionen der Planeten angibt, die sich notfalls ebenso leicht aus den geozentrischen Positionen ableiten lassen. Wir könnten hier noch manche anderen Wünsche anfügen, zu denen uns der langjährige Gebrauch astronomischer Ephemeriden veranlasst hat. Da sie sich jedoch kaum ohne größere Änderungen an deren Form und Einrichtung erfüllen ließen, beschränken wir uns hier auf einen einzigen Wunsch. Wenn der verdienstvolle Herausgeber ihn erfüllte, würde er allen einen großen Gefallen tun, die viel mit Planeten- und Kometenrechnungen zu tun haben und für die das Berechnen der Sonnenpositionen stets sehr viel Zeit beansprucht. Dieser Zeitaufwand ließe sich in sehr vielen Fällen ganz vermeiden, wenn die Abstände der Sonne von der Erde oder ihre Logarithmen für jeden einzelnen Tag des Jahres angegeben würden, notfalls sogar mit einer Dezimalstelle weniger, als Herr \textsc{Bode} in den letzten vier Jahrgängen angegeben hat. Die dadurch etwas vermehrte Arbeit könnte der Herausgeber ausgleichen, indem er bei den Längen und Breiten des Mondes die Sekunden wegließe. Ihre Berechnung kostet ungeheure Arbeit, und wir glauben behaupten zu können, dass sie kaum oder überhaupt nicht gebraucht werden.

Zu diesen Bemerkungen über die Einrichtung des Jahrbuchs hat uns allein der aufrichtige Wunsch bewogen, es noch brauchbarer zu sehen. Nun wenden wir uns den Abhandlungen zu, die dem astronomischen Jahrbuch bleibenden Wert geben. 1) Untersuchung der wahren elliptischen Bewegung des Kometen von 1769, von \textsc{Fr. W. Bessel}. Diese Schrift wurde für einen Wettbewerb um einen Preis von 30 Friedrichsd’or eingereicht, den ein Ungenannter für die beste astronomische Abhandlung oder die bemerkenswerteste Entdeckung ausgesetzt hatte. Man sprach Herrn \textsc{Bessel} jedoch nur die eine Hälfte zu und die andere Herrn \textsc{Huth} als Mitentdecker der beiden Kometen des Jahres 1805. Die elliptische Bahn jenes Kometen hatte schon mehrere Astronomen beschäftigt, unter denen Pater \textsc{Asklepi} sie am sorgfältigsten untersucht hatte. Der Gegenstand verdiente aber durchaus eine neue Bearbeitung. Dazu luden vor allem die großen Verbesserungen unserer Sternverzeichnisse und Sonnentafeln ein, mit deren Hilfe man auf viel zuverlässigere Ergebnisse hoffen durfte. Herr \textsc{Bessel} hat ihn jetzt so vorbildlich und erschöpfend behandelt, dass nichts zu wünschen übrig bleibt. Wenn man zugestehen darf, dass die von Herrn \textsc{Bessel} aus einer neuen, äußerst sorgfältigen Auswertung sämtlicher guter Beobachtungen abgeleiteten Fundamentalpositionen bis auf $5''$ in Rektaszension und Deklination zuverlässig sind, so liegt die Umlaufzeit des Kometen zwischen $1691\frac{3}{5}$ und 2673 Jahren. Besondere Aufmerksamkeit verdient das Verfahren von Herrn \textsc{Bessel} zur Berechnung der Wirkung der Refraktion. Nur schade, dass die dazugehörigen Formeln durch mehrere Druckfehler völlig entstellt sind und man sie deshalb ohne einen späteren ausführlichen Aufsatz über diesen Gegenstand in der Mon. Corresp. gar nicht gebrauchen kann. 2) Über die geografische Länge von Havanna, von Herrn \textsc{Oltmanns}. Nach Beobachtungen von \textsc{Churrucca}, \textsc{Robredo}, \textsc{v. Humboldt} und \textsc{Galiano} beträgt sie im Mittel $5^{\mathrm u}\,38^{\mathrm m}\,51^{\mathrm s}5$ westlich von Paris. 3) Tafeln zur Reduktion der Höhen des Polarsterns auf die Meridianhöhe, für die Breite von Berlin. 4) Scheinbare Helligkeitsänderungen von Algol in den Jahren 1808, 1809, 1810, von Herrn Professor \textsc{Wurm}. 5) Drittes und viertes Verzeichnis der verglichenen Helligkeiten der Fixsterne, von Herrn \textsc{Herschel}. 6) Über die wahre geografische Länge des in Peru gemessenen Breitengrads, von Herrn \textsc{Oltmanns}. Aus den Beobachtungen von \textsc{Humboldt}, \textsc{Ulloa} und den französischen Akademikern ergibt sich als mittlere Länge von Quito $5^{\mathrm u}\,24^{\mathrm m}\,20^{\mathrm s}$. Der letzte Band der Connoissance des tems für 1809 gab dafür nach \textsc{Bouguer} noch $5^{\mathrm u}\,21^{\mathrm m}\,0^{\mathrm s}$ an. 7) Tafeln zur Berechnung der jährlichen Änderung der Rektaszension und Deklination der Fixsterne. In diesen Aufsatz haben sich mehrere Unrichtigkeiten eingeschlichen. Die Wirkung von Sonne und Mond bringt nicht die jähr\SourceLeaf{511}liche Änderung von $50''15$ in der Länge hervor, die sich durch die Änderung der Schiefe der Ekliptik noch um $0''135$ vergrößert, wie der Verfasser des Aufsatzes es ausdrückt. Vielmehr wird die durch die Wirkung von Sonne und Mond erzeugte sogenannte lunisolare Präzession durch die Verschiebung der Ekliptik wieder etwas vermindert. Diese Verschiebung hat zwar dieselbe Ursache wie die Abnahme der Schiefe, darf aber nicht mit ihr verwechselt werden. Dadurch sinkt die Präzession auf $50''15$ oder nach genaueren Bestimmungen auf $50''10$, den beobachteten Wert. Die Koeffizienten des veränderlichen Teils der Präzession in Rektaszension und die Formel für die Änderung der Deklination sollten ganz gleich sein. Auch die Regel zur Berechnung der Änderung der jährlichen Variation ist falsch. 8) Beitrag zu den Verfahren, eine Reihe von Monddistanzen zur Berechnung geografischer Längen heranzuziehen, von Herrn \textsc{Jabbo Oltmanns} (der Unterschied zwischen scheinbarer und wahrer Distanz wird für drei Zeitpunkte berechnet und daraus durch Interpolation für alle Beobachtungen bestimmt). 9) Astronomische Beobachtungen in Prag 1806, von den Herren \textsc{David} und \textsc{Bittner}. 10) Methode zur Bestimmung der geografischen Position von Gegenständen auf der Erde mithilfe beobachteter Azimute, Höhenwinkel und relativer Höhen, nebst einigen Hilfstafeln zur Berechnung barometrischer Höhenmessungen, von Herrn \textsc{Jabbo Oltmanns}. Viel Genauigkeit kann dieses Verfahren freilich nicht liefern. In ganz unbekannten Gegenden, wo andere Beobachtungen fehlen, ist es aber nicht zu verwerfen. 11) Über das röhrenförmige Pendel von \textsc{Troughton}, von Herrn \textsc{Schnitter} in Aachen. Im Jahrbuch für 1808 war die Erfindung von \textsc{Troughton} so beschrieben, dass sich selbst mit den Angaben des Erfinders für die Ausdehnung von Messing und Stahl nicht die richtige Kompensation ergab. Herrn \textsc{Schnitter} gebührt Dank dafür, dass er darauf aufmerksam gemacht hat und darauf hinwies, dass für eine vollkommene Kompensation alle sich nach unten ausdehnenden Stangen aus Stahl bestehen müssen. Der Rezensent, der den Originalaufsatz von \textsc{Troughton} (\textsc{Nicholson}’s philosophical journal, Bd. IX) vor sich hat, glaubt allerdings darauf hinweisen zu müssen, dass der Irrtum allein in der falschen Übersetzung im Jahrbuch liegt. Die Vorschrift von \textsc{Troughton} ist genau so, wie sie nach den Angaben von Herrn \textsc{Schnitter} sein muss. 12) Astronomische Beobachtungen und Bemerkungen von Herrn \textsc{van Beeck Calkoen}. 13) Beobachtung der Bedeckung von $\alpha$ im Skorpion am 2c.\ März 1805 und der Sonnenfinsternis vom 16. Juni 1806 auf der Insel Leon bei Cadiz, von Herrn \textsc{Canelas}. 14) Astronomische Beobachtungen in Wien 1806, von Herrn \textsc{Triesnecker}. 15) Entdeckung und Beobachtung eines vierten neuen Planeten von Herrn Dr. \textsc{Olbers}. 16) Beobachtung und Berechnung der Bahn des Kometen von 1806, Beobachtungen von Vesta und Juno und von Sternbedeckungen, von Herrn F. W. \textsc{Bessel}. 17) Messung der scheinbaren Größe der Vesta, von Herrn \textsc{Schröter} (am 26. April 1807 betrug der scheinbare Durchmesser $0''488$, woraus ein wahrer Durchmesser von 68 Meilen folgt). 18) Beobachtung der Vesta und Berechnung der Elemente ihrer Bahn und ihres Laufs für 1808, von Herrn Dr. \textsc{Gauss}. 19) Beobachtungen der Vesta in Berlin, vom Herausgeber. 20) Astronomische Beobachtungen und Bemerkungen von Herrn \textsc{Fritsch} in Quedlinburg. Die Beobachtungen der Planeten, auch die der Juno im Jahr 1807, als sie sehr schwer zu beobachten war, machen Herrn \textsc{Fritsch} Ehre, auch wenn sie für den Astronomen mangels einer gehörigen Zeitbestimmung nicht brauchbar sind. 21) Über ein Quecksilberpendel, von Herrn \textsc{Th. Blacker} in London. 22) Astronomische Beobachtungen von Herrn \textsc{Lalande}. 23) Beobachtungen zu den klimatischen Verhältnissen und zur Atmosphäre des Saturn, von Herrn Dr. \textsc{Herschel}. Die von Herrn Dr. H. bemerkten Farbänderungen der Polargegenden des Saturn scheinen periodisch zu sein. Herr H. sieht sie als Beweis für das Vorhandensein einer Atmosphäre an. Astronomische Beobachtungen an der Berliner Sternwarte 1806. Die Kriegsunruhen im Herbst dieses Jahres führten dabei zu mancherlei Störungen und Unterbrechungen. 25) Lauf der Pallas und Juno im Jahr 1808 (vorausberechnet). 26) Über beobachtete Unterschiede in den scheinbaren Größen einiger Sterne, von Herrn Dr. \textsc{Koch} in Danzig. 27) Astronomische Beobachtun\SourceLeaf{512}gen in Spanien (einigen spanischen Schriften entnommen). 28) Einige Bemerkungen zur physikalischen Astronomie, von Herrn \textsc{Huth} in Frankfurt an der Oder. 29) Beitrag zu geografischen Längenbestimmungen, von Herrn \textsc{Oltmanns}. 30) Der Vorschlag des Grafen \textsc{von Platen} für ein Verfahren zur Lösung einer astronomischen Aufgabe beruht im Wesentlichen darauf, Produkte von Sinuswerten in Summen umzuwandeln, wie es die Astronomen vor der Erfindung der Logarithmen taten. Heute rechnet man in den meisten Fällen bequemer mit Logarithmen. Gerade in den Fällen, von denen der Verfasser spricht, findet der Rezensent das vorgeschlagene Verfahren nicht bequem. 31) Vorschlag für ein Verfahren zur Bestimmung der Horizontalrefraktion mithilfe der geografischen Länge, von Herrn \textsc{Oltmanns} (mithilfe der Abstände des Mondes von der tief stehenden Sonne). 32) Zusatz von Herrn \textsc{Fischer} zu der Abhandlung über die beste Form der Objektivspiegel im Jahrbuch 1806. Herr F. nimmt hier seine irrige Behauptung zurück, dass sphärische Spiegel den parabolischen vorzuziehen seien. 33) Verschiedene astronomische Beobachtungen und Nachrichten. Aufmerksamkeit verdient die Mitteilung über einen Stern fünfter Größe, den Herr \textsc{Hecker} in Remplin bei $\alpha$ im Füllen am 17. August 1804 beobachtete und der danach nicht wieder gesehen wurde. Wir wundern uns nur, dass Herr \textsc{Hecker} dies nicht sogleich mitgeteilt hat, als die Nachricht vielleicht von Nutzen hätte sein können.
